\section{Background}\label{background}

This section formalizes the exact fully dynamic kNN join problem and reviews the two index structures that serve as the main building blocks of \dualtree.

\subsection{Definitions and PCA}\label{sec:background-preliminaries}

Let $\mathcal{X}\subseteq\mathbb{R}^D$ be a $D$-dimensional vector space. Unless stated otherwise, $dist(\cdot,\cdot)$ denotes Euclidean distance, i.e., the $\ell_2$ norm.

\begin{definition}[kNN search]\label{def:knns}
Given a query point $u\in\mathcal{X}$, a reference set $I\subseteq\mathcal{X}$, and an integer $k\in\{1,\ldots,|I|\}$, $\knn(u,I)$ denotes the set of $k$ points in $I$ with smallest distance to $u$. Ties are resolved by a fixed total order over point identifiers.
\end{definition}

\begin{definition}[kNN distance]\label{def:dknn}
Given $u$, $I$, and $k$, the kNN distance of $u$ over $I$ is the distance from $u$ to the farthest point in its kNN result:
\[
\dknn(u,I)=\max_{x\in \knn(u,I)} dist(u,x).
\]
\end{definition}

\begin{definition}[RkNN search]\label{def:rknn}
Given a query set $U\subseteq\mathcal{X}$, a reference set $I\subseteq\mathcal{X}$, an integer $k$, and a reference point $i\in I$, the RkNN result of $i$ is
\[
\rknn(i,U,I)=\{u\in U \mid i\in \knn(u,I)\}.
\]
Equivalently, under the tie-breaking rule above, $u$ is returned when $dist(u,i)\le \dknn(u,I)$ and $i$ is admitted by the induced ordering.
\end{definition}

\begin{definition}[kNN join]\label{def:knnj}
Given $U\subseteq\mathcal{X}$, $I\subseteq\mathcal{X}$, and $k$, the kNN join from $U$ to $I$ is the materialized table
\begin{equation}
  \knnjoin(U,I)=
  \{(u,\knn(u,I))\mid u\in U\}.
\end{equation}
\end{definition}

Principal component analysis (PCA) is a standard dimensionality-reduction technique~\cite{wold1987principal,chakrabarti2000local,abdi2010principal,vidal2016principal}. We use PCA for exact pruning: projected distances in a prefix of principal components provide lower bounds on full-dimensional Euclidean distances. Let $PC(i)$ denote the $i$-th principal component, and let $\cup_{i=1}^{d}PC(i)$ denote the $d$-dimensional subspace spanned by the leading $d$ components. For a point $u$, $u_{\cup_{i=1}^{d}PC(i)}$ denotes its projection into this subspace.

\begin{lemma}[Contractive distance property of PCA]\label{lm:pca_decrease_dist}\label{thm:distance_pca}
For any two points $a,b\in\mathcal{X}$ and any $d_1\le d_2\le D$,
\[
dist_{\cup_{i=1}^{d_1}PC(i)}(a,b)
\le
dist_{\cup_{i=1}^{d_2}PC(i)}(a,b)
\le
dist(a,b).
\]
\end{lemma}

This monotonicity follows from Euclidean distance being a sum of nonnegative squared coordinate differences; adding projected dimensions cannot decrease the partial sum. A proof of this property was also given in the earlier work~\cite{ukey2023efficient}. In our setting, this fact is the correctness basis for all reduced-space pruning steps used later in the paper, including the cluster-level pruning in \deltatree and \hdrtree and the point-wise pruning stage introduced in Section~\ref{sec:pruning}.

\subsection{Problem Statement}\label{sec:problem-statement}

We study exact fully dynamic kNN join in high dimensions: exactly maintain $\knnjoin(U,I)$ efficiently as both $U$ and $I$ are updated, with $D$ large enough that full-dimensional distance computation is a heavy cost.

\subsection{Index Building Blocks}\label{sec:index-building-blocks}

\deltatree~\cite{cui2003contorting} and \hdrtree~\cite{yang2014continuous} are two closely related high-dimensional hierarchical indexes. They share a similar top-down organization and rely on reduced-space pruning, but they are designed for different search tasks. This combination of hierarchical clustering, dimensionality reduction, and exact leaf-level verification is precisely what makes them effective building blocks for exact search in high-dimensional spaces. \deltatree indexes the reference set $I$ and accelerates kNN search, whereas \hdrtree indexes the query set $U$ and accelerates RkNN search. In our setting, this means \deltatree supports the kNN computations needed to create or repair join rows, whereas \hdrtree supports the identification of affected queries when reference updates occur. We therefore describe their common structure first and then highlight their task-specific roles.

\noindent\textbf{\textit{Structure.}}
Both indexes partition points top down. At level $l$, points are represented in a PCA subspace of dimensionality $d_l$, where the dimensionality increases with depth. A non-leaf node stores cluster summaries including the cluster center, radius, level information, and child pointers. If a cluster contains more than a threshold number of points, it is recursively partitioned; otherwise it becomes a leaf that stores the original full-dimensional points. This shared organization supports inexpensive coarse pruning near the root and postpones exact full-dimensional verification until the leaf level.

\noindent\textbf{\textit{Search.}}
The two indexes differ in the search primitive they support. For kNN search on \deltatree, clusters whose low-dimensional lower bound already exceeds the current $k$-th best distance can be discarded. For RkNN search on \hdrtree, clusters whose low-dimensional lower bound exceeds the maximum stored kNN distance of their members can be discarded. In later sections we denote this cluster-level quantity by $\dmknn(C)$, namely the maximum $\dknn(u,I)$ over the queries stored in the subtree rooted at cluster $C$. In both cases, pruning is driven by reduced-space lower bounds and surviving candidates are verified in the original space, with correctness following from Lemma~\ref{lm:pca_decrease_dist}. This difference in search semantics is exactly why \deltatree fits row computation and repair through kNN search, whereas \hdrtree fits affected-query identification through RkNN search.

\noindent\textbf{\textit{Maintenance of \deltatree and \hdrtree.}}
\deltatree provides update routines for insertions and deletions in $I$~\cite{cui2003contorting}. \hdrtree, in contrast, was introduced for a static query set and a dynamic reference set, so query-side insertions and deletions are not directly covered by its original setting. In our framework, we extend \hdrtree into a dynamically maintainable structure for affected-query localization so that updates on the query side can be absorbed as first-class operations. This gives \dualtree the movable components required on both sides. On top of that, the framework coordinates these components with the materialized join table and their stored index summaries: updates on one side may trigger repairs of join rows, and those repairs in turn may require synchronized updates to the pruning values maintained in the two indexes. Section~\ref{sec:framework} explains how this coordinated maintenance loop is organized under different update cases.

Taken together, these properties explain why \deltatree and \hdrtree are used as the foundations of \dualtree. Our setting is exact fully dynamic maintenance in a high-dimensional space. This requires the index structures must support efficient pruning under high dimensionality, preserve exact verification, and remain maintainable under repeated updates. To the best of our knowledge, and consistent with the survey of exact high-dimensional search methods~\cite{ukey2023survey}, \deltatree and \hdrtree remain state-of-the-art exact structures for the two search tasks required here: high-dimensional kNN search for row computation and repair, and high-dimensional RkNN search for locating affected queries. Many other methods satisfy only part of this requirement. Mapping-based approaches such as iDistance~\cite{jagadish2005idistance} and standalone exact filtering methods such as PL-Tree~\cite{wang2013pl} or triangle-inequality-based search~\cite{pan2020new} are effective for exact kNN retrieval itself, but they are less suitable as components of our framework because they do not provide reduced-space hierarchical pruning. Join methods such as MuX~\cite{bohm2001cost,bohm2004k}, Gorder~\cite{xia2004gorder}, iJoin~\cite{yu2007efficient}, and Sphere tree ~\cite{yu2010high} are also not suitable as the core search structures of our framework, because they are either difficult to extend into repeated fully dynamic maintenance components or are not targeted at the exact high-dimensional search regime required here. Faster recent alternatives are predominantly approximate, which changes the contract of the problem rather than solving the exact maintenance task studied in this paper. By contrast, \deltatree and \hdrtree remain particularly strong exact building blocks for our problem.
